\documentclass[10pt]{article}
\usepackage[margin=0.8in]{geometry}
\usepackage{amsmath}
\usepackage{booktabs}
\usepackage{float}
\title{A state-space increment with an Armijo line search on the GridFM-PF backbone\\
\large A negative result, and what it does and does not show}
\author{PIGNN-Attn-LS-PPC experimental summary}
\date{14 August 2026}
\begin{document}
\maketitle

\section{Why this was tried}

Reading GENCO (arXiv:2608.09921) raised a question: GENCO runs $N$ correction
steps, PIGNN-Attn-LS runs $K=40$ correction steps, so if a line search helps the
latter it should help the former. Examining the code showed the premise was only
half right.

The released GridFM-PF model (\texttt{GNS\_heterogeneous}) already runs the
iterative loop with physics feedback. Its per-layer body decodes an
\emph{absolute} voltage state from the latent, computes the power-balance
residual, and injects it back:

\begin{center}
\texttt{h\_bus = h\_bus + physics\_mlp(bus\_residuals)}
\end{center}

GENCO's Eq.~(17) is the same update. So the feedback loop is not GENCO's
contribution over GridFM-PF --- its contributions are the HGT backbone, the
discounted deep-supervision loss, the OPF/SE decoders, box constraints, and
analytic $Q_g$ recovery.

\paragraph{The structural difference that matters.} The GridFM/GENCO update adds
in \emph{latent} space and re-decodes the state from scratch at every layer.
PIGNN updates the \emph{state}, $V \leftarrow \mathrm{clamp}(V + \alpha\,dV)$,
with $\alpha$ from Armijo backtracking. A line search needs a step in the space
the objective is a function of; a latent injection is not one. So a line search
cannot be attached to GridFM-PF or GENCO as they stand.

\paragraph{The change.} Read the per-layer decoder output as an increment to a
running state instead of as a fresh absolute state, and scale it by a line
search on the power-balance residual. Everything else --- backbone, feedback
loop, masking, loss, data, splits --- unchanged. Three variants were run:

\begin{table}[H]
\centering\small
\begin{tabular}{llll}
\toprule
Tag & Decoder output & Step length & Parameters \\
\midrule
\texttt{graphkit}     & absolute state (upstream) & --- & 20{,}069{,}187 \\
\texttt{graphkit\_inc} & increment                & always 1 & 20{,}069{,}187 \\
\texttt{graphkit\_ls}  & increment                & Armijo backtracking & 20{,}069{,}187 \\
\bottomrule
\end{tabular}
\caption{Identical parameter counts, confirmed at run time: the backbone was not
touched, only the forward pass. \texttt{graphkit\_inc} exists to separate the
reparameterisation from the line search --- without it a result could be
attributed to either.}
\end{table}

\section{Result: the change makes the model worse}

\texttt{graphkit\_ls} is worse than the unmodified \texttt{graphkit} on
\textbf{22 of the 26} grids that finished, median $|V|$ RMSE ratio
\textbf{1.81}, and worse on $\Delta P^{\infty}$ on 23 of 26, median ratio
\textbf{3.75}.

\begin{table}[H]
\centering\scriptsize
\begin{tabular}{lrrrrrrr}
\toprule
 & & \multicolumn{3}{c}{$|V|$ RMSE [pu]} & \multicolumn{3}{c}{$\Delta P^{\infty}$ [pu]} \\
\cmidrule(lr){3-5}\cmidrule(lr){6-8}
Grid & Buses & graphkit & \_ls & ratio & graphkit & \_ls & ratio \\
\midrule
case5 & 5 & 6.459e-05 & 1.241e-02 & 192.1 & 3.82e-02 & 1.35e+01 & 354.0 \\
case6ww & 6 & 4.757e-04 & 4.674e-02 & 98.3 & 1.57e-02 & 1.51e+00 & 95.8 \\
case14 & 14 & 8.409e-04 & 2.563e-02 & 30.5 & 2.65e-02 & 2.61e+00 & 98.3 \\
SimBench & 94 & 1.584e-03 & 3.793e-02 & 23.9 & 8.74e-02 & 2.61e+00 & 29.8 \\
case1888rte & 1888 & 2.887e-02 & 2.119e-01 & 7.3 & 1.38e+02 & 8.82e+02 & 6.4 \\
case89pegase & 89 & 1.285e-02 & 5.995e-02 & 4.7 & 1.43e+00 & 2.13e+02 & 148.6 \\
case9 & 9 & 8.752e-04 & 3.327e-03 & 3.8 & 2.93e-02 & 1.01e-01 & 3.4 \\
GBreducednetwork & 29 & 5.086e-04 & 1.566e-03 & 3.1 & 1.52e+01 & 1.91e+01 & 1.3 \\
LVN\_heo1 & 722 & 2.229e-02 & 5.100e-02 & 2.3 & 6.51e+00 & 1.76e+02 & 27.1 \\
case4gs & 4 & 2.857e-04 & 6.244e-04 & 2.2 & 2.05e-02 & 6.69e-02 & 3.3 \\
case30 & 30 & 1.753e-03 & 3.647e-03 & 2.1 & 3.44e-02 & 1.03e-01 & 3.0 \\
case2848rte & 2848 & 3.176e-02 & 6.335e-02 & 2.0 & 2.72e+01 & 2.86e+03 & 105.2 \\
case\_ieee30 & 30 & 4.270e-03 & 7.800e-03 & 1.8 & 5.91e-02 & 1.92e-01 & 3.2 \\
iceland & 189 & 1.586e-02 & 2.852e-02 & 1.8 & 1.15e+00 & 4.37e+02 & 380.1 \\
case57 & 57 & 6.306e-03 & 1.090e-02 & 1.7 & 1.49e-01 & 1.83e-01 & 1.2 \\
case33bw & 33 & 9.452e-03 & 1.621e-02 & 1.7 & 2.06e-02 & 8.97e-02 & 4.3 \\
case118 & 118 & 1.320e-03 & 2.246e-03 & 1.7 & 2.61e-01 & 2.71e-01 & 1.0 \\
GBnetwork & 2224 & 3.354e-02 & 5.656e-02 & 1.7 & 1.70e+02 & 1.51e+04 & 88.5 \\
case2869pegase & 2869 & 2.860e-02 & 3.325e-02 & 1.2 & 1.92e+02 & 1.49e+03 & 7.8 \\
case\_illinois200 & 200 & 7.814e-03 & 8.444e-03 & 1.1 & 2.60e-01 & 5.18e-01 & 2.0 \\
case3120sp & 3120 & 2.771e-02 & 2.834e-02 & 1.0 & 3.28e+03 & 1.62e+03 & 0.5 \\
case300 & 300 & 2.343e-02 & 2.343e-02 & 1.0 & 1.03e+01 & 3.52e+01 & 3.4 \\
case1354pegase & 1354 & 2.324e-02 & 1.995e-02 & 0.9 & 2.90e+01 & 1.18e+02 & 4.1 \\
case24\_ieee\_rts & 24 & 1.677e-03 & 1.289e-03 & 0.8 & 2.32e-01 & 1.73e-01 & 0.7 \\
case39 & 39 & 9.046e-03 & 6.804e-03 & 0.8 & 6.92e-01 & 7.52e-01 & 1.1 \\
case145 & 145 & 1.666e-01 & 2.792e-02 & 0.2 & 3.43e+02 & 2.97e+01 & 0.1 \\
\bottomrule
\end{tabular}
\caption{Sorted worst-first by $|V|$ ratio. Ratios above 1 mean the modified
model is worse.}
\end{table}

\section{It is not a large-grid win}

The four grids where \texttt{\_ls} wins invite the reading that the change helps
on large grids. It does not survive checking.

\begin{table}[H]
\centering\small
\begin{tabular}{lrrr}
\toprule
 & $n$ & median $|V|$ ratio & \_ls better on \\
\midrule
$<1000$ buses    & 20 & 1.95 & 3 \\
$\geq 1000$ buses & 6 & 1.42 & 1 \\
\bottomrule
\end{tabular}
\caption{Both medians are above 1: on large grids the modified model is less bad,
not better.}
\end{table}

Two facts settle it. First, the individually largest grids are among the worst
losses --- case1888rte $7.3\times$, case2848rte $2.0\times$, GBnetwork
$1.7\times$. Second, the rank correlation of $\log$(ratio) is
$-0.451$ against $\log$(bus count) but \textbf{$-0.539$ against the baseline's
own $|V|$ RMSE}: the pattern tracks how good the baseline already was more
closely than it tracks size, and size and baseline difficulty are themselves
correlated.

The wins are consistent with a ceiling effect. Two of the four are small grids
(case24\_ieee\_rts at 24 buses, case39 at 39). The largest win, case145 at
$0.17\times$, is a grid where the baseline is \emph{broken}: graphkit posts a
negative $|V|$ regression slope there ($-0.199$ on PQ buses), i.e.\ its
prediction is anti-correlated with the reference. Beating a broken baseline is
not evidence of benefit.

\section{Which half of the change is at fault}

This is the part worth keeping. On the 11 grids where all three variants
finished:

\begin{table}[H]
\centering\scriptsize
\begin{tabular}{lrrrrrrr}
\toprule
 & & \multicolumn{3}{c}{$|V|$ RMSE [pu]} & \multicolumn{3}{c}{ratio} \\
\cmidrule(lr){3-5}\cmidrule(lr){6-8}
Grid & Buses & graphkit & \_inc & \_ls & \_inc/gk & \_ls/gk & \_ls/\_inc \\
\midrule
case4gs & 4 & 2.857e-04 & 3.670e-03 & 6.244e-04 & 12.84 & 2.19 & 0.17 \\
case5 & 5 & 6.459e-05 & 6.867e-04 & 1.241e-02 & 10.63 & 192.11 & 18.07 \\
case6ww & 6 & 4.757e-04 & 3.954e-03 & 4.674e-02 & 8.31 & 98.26 & 11.82 \\
case9 & 9 & 8.752e-04 & 3.125e-03 & 3.327e-03 & 3.57 & 3.80 & 1.06 \\
case14 & 14 & 8.409e-04 & 8.489e-03 & 2.563e-02 & 10.10 & 30.48 & 3.02 \\
case24\_ieee\_rts & 24 & 1.677e-03 & 1.092e-02 & 1.289e-03 & 6.51 & 0.77 & 0.12 \\
GBreducednetwork & 29 & 5.086e-04 & 6.221e-03 & 1.566e-03 & 12.23 & 3.08 & 0.25 \\
case30 & 30 & 1.753e-03 & 7.993e-03 & 3.647e-03 & 4.56 & 2.08 & 0.46 \\
case\_ieee30 & 30 & 4.270e-03 & 1.451e-02 & 7.800e-03 & 3.40 & 1.83 & 0.54 \\
case39 & 39 & 9.046e-03 & 8.047e-03 & 6.804e-03 & 0.89 & 0.75 & 0.85 \\
case145 & 145 & 1.666e-01 & 3.670e-02 & 2.792e-02 & 0.22 & 0.17 & 0.76 \\
\bottomrule
\end{tabular}
\end{table}

\begin{table}[H]
\centering\small
\begin{tabular}{lrr}
\toprule
Comparison & median ratio & better on \\
\midrule
\texttt{\_inc} vs \texttt{graphkit} & 6.51 & 2/11 \\
\texttt{\_ls} vs \texttt{graphkit}  & 2.19 & 3/11 \\
\textbf{\texttt{\_ls} vs \texttt{\_inc}} & \textbf{0.76} & \textbf{7/11} \\
\bottomrule
\end{tabular}
\end{table}

\textbf{The incremental reparameterisation is what breaks the model, and the line
search partially repairs it.} Turning the decoder into an increment costs a
factor of 6.5; adding the line search on top recovers roughly a third of that
($0.76\times$, better on 7 of 11). This inverts the hypothesis the experiment
started from --- the working assumption before running \texttt{\_inc} was that
the line search was suppressing steps and causing the damage. It is doing the
opposite.

\paragraph{Why the increment hurts.} The upstream decoder is trained to emit an
absolute state and is followed immediately by a physics decoder that completes
the solution. Re-reading its output as a correction to a running state changes
what the head must learn without giving it any more capacity, and it removes the
model's ability to discard a bad intermediate state --- an absolute decode can
jump anywhere, an increment can only walk.

\section{Corroborating measurement: what the line search actually did}

A probe of a smoke-trained checkpoint, 12 layers $\times$ 32 graphs:

\begin{table}[H]
\centering\small
\begin{tabular}{lr}
\toprule
mean $\alpha$ & 0.115 \\
fraction $\alpha = 1$ & 0.000 \\
fraction $\alpha = 0$ (step rejected outright) & 0.354 \\
layers with $\alpha = 0$ on every graph & 3 of 12 \\
\bottomrule
\end{tabular}
\end{table}

The full step was never accepted and a third of steps were rejected, which
prompted widening the backtracking range from 4 to 6 halvings before the
campaign. In hindsight this was treating a symptom: the acceptance test is
applied to a direction with no descent guarantee, because it comes from a
learned head rather than a gradient. Rejection therefore does not mean "the step
was too long", and lowering the floor cannot fix it.

\section{What this does and does not show}

\paragraph{It does not contradict the line search helping PIGNN-Attn-LS.} PIGNN
constructs its direction to be a descent direction --- it projects onto the
descent half-space before backtracking --- and its increments are the model's
native parameterisation, not a reinterpretation of a head trained for something
else. Nothing here bears on that.

\paragraph{It does show that the ingredient is not portable by itself.} Bolting a
state-space line search onto a backbone whose decoder was trained to emit
absolute states does not transfer the benefit; the reparameterisation it requires
costs more than the line search returns.

\paragraph{What a fair test would need.} Train the incremental decoder from a
matched initialisation with a descent-projected direction, as PIGNN does, rather
than reinterpreting the upstream head. That is a different and larger experiment
than this one, and it was not run.

\section{Status and coverage}

The experiment is \textbf{abandoned}; all jobs on both clusters were cancelled.
Nothing here is a complete campaign and none of it should be folded into the
31-grid comparison tables.

\begin{table}[H]
\centering\small
\begin{tabular}{llrl}
\toprule
Variant & Cluster & Finished & Notes \\
\midrule
\texttt{graphkit\_ls}  & helma h100 & 26/31 & 1 NODE\_FAIL (case6515rte), 4 cancelled mid-run \\
\texttt{graphkit\_inc} & alex2 a100 & 11/31 & 20 cancelled; the 11 are all $\leq 145$ buses \\
\bottomrule
\end{tabular}
\end{table}

Two caveats on the numbers above. The three-way decomposition rests on 11 small
grids only, because \texttt{\_inc} was cancelled before any large grid finished
--- the direction of that result is clear but its magnitude on large grids is
unmeasured. And \texttt{\_inc} ran on a100 with a smaller batch budget
(\texttt{BN\_BUDGET} 50000 against 120000) because an a100 is a 40~GB card; the
batch cap of 64 binds below roughly 780 buses, so for all 11 grids in the
decomposition the batch was identical to the helma runs and that confound does
not apply here.

\paragraph{Code.} \texttt{genco\_ls.py} and the
\texttt{-{}-gridfm\_impl graphkit\_ls|graphkit\_inc} switches are kept. They are
a working, parameter-identical implementation of the idea and are the starting
point if the fair test above is ever run.

\end{document}
